% ==================== PREÁMBULO ====================
% Archivo para centralizar la configuración de paquetes y estilos del documento.

% ----- IDIOMA Y CODIFICACIÓN -----
\usepackage[T1]{fontenc}                % Codificación de fuentes para caracteres especiales.
\usepackage[utf8]{inputenc}             % Codificación de entrada para acentos y eñes.
\usepackage[spanish,english]{babel}

%\usepackage[spanish,es-nosectiondot,es-lcroman]{babel} % Soporte para español.

% ----- PAQUETES MATEMÁTICOS -----
\usepackage{amsmath,amssymb,amsfonts}   % Símbolos y entornos matemáticos de AMS.

% ----- FIGURAS Y GRÁFICOS -----
\usepackage{graphicx}                   % Para incluir imágenes (jpg, png, etc.).
\usepackage{xcolor}                     % Para definir y usar colores.
\usepackage{setspace}                   % Para dejar espacio.
\usepackage{enumitem}
\usepackage{titlesec}

% ----- TABLAS -----
\usepackage{booktabs}                   % Líneas de calidad profesional en tablas.
\usepackage{tabularx}                   % Tablas con columnas de ancho ajustable.
\usepackage{multirow}                   % Celdas que ocupan múltiples filas.
\usepackage{multicol}                   % Para usar múltiples columnas en el texto.
\usepackage{threeparttable}             % Permite añadir notas al pie dentro de las tablas.
\usepackage{float}                      % Mejora el control sobre la posición de flotantes (figuras, tablas).
\usepackage{booktabs} % Para líneas más elegantes

% ----- CODIGOS -----
% \usepackage{algorithm}
\usepackage{algpseudocode}
\usepackage{amsmath} % Imprescindible para \text{}
\usepackage{listings}

% ----- BIBLIOGRAFÍA Y REFERENCIAS -----
% \usepackage[numbers]{natbib}            % Gestión de bibliografía estilo autor-año o numérico.
\usepackage{hyperref}                   % Creación de hipervínculos y referencias cruzadas.
\hypersetup{
  colorlinks=true,                    % Colorea los enlaces en lugar de usar recuadros.
  breaklinks=true                     % Permite que los enlaces largos se dividan en varias líneas.
}
\usepackage{lineno}                     % Para numerar las líneas del documento.

% ----- ALGORITMOS -----
\usepackage{algorithmic}                % Para escribir pseudocódigo.

% ----- AJUSTES TIPOGRÁFICOS -----
\usepackage{etoolbox}                   % Herramientas de programación para LaTeX (usado para bibliografía).
\hbadness=10000                         % Suprime advertencias de "underfull \hbox".
\vbadness=10000                         % Suprime advertencias de "underfull \vbox".
\tolerance=10000                        % Aumenta la tolerancia para evitar líneas que sobresalen del margen.

% ==================== CONSEJOS RÁPIDOS ====================
% - Citas: \cite{clave} o \citep{clave} con natbib.
% - Figuras: \includegraphics[width=...]{ruta/figura.png} o entornos TikZ/PGF.
% - Tablas: \begin{table} ... \end{table}, usando booktabs para mejor estilo.
% - Compilación: Se recomienda usar latexmk o la secuencia pdflatex -> bibtex -> pdflatex -> pdflatex.
